\section{Comparison along uniformly allocated orbit blocks}
\label{sec:dynamical-comparison}

We now obtain the bundle embedding needed in
\Cref{thm:torus-compression} from a comparison assumption in the crossed
product.  There are two losses to control.  Only a small fraction of orbit
blocks fail to have a continuous choice of alignment, and only a small
fraction of the sparse parameter slots are nonzero.  The corresponding
bundle summands will be embedded separately.  All other summands are
handled simultaneously by a single Cuntz comparison witness.

Throughout this section, $h:X\to X$ is a minimal homeomorphism of an
infinite compact metrizable space and $A=C(X)\rtimes_h\mathbb Z$.
Write $\mathcal M_h(X)$ for the invariant Borel probability measures.
We use comparison bounds in the sense of \Cref{eq:comparison-bound},
with normalized traces on $A$ and unnormalized matrix traces.

\subsection{Open towers and uniformly good starting points}

The dynamical input is the following special case of
\cite[Theorem~4.2]{Niu2024}. Niu obtains this tower from the continuous
equivariant tilings of Gutman, Lindenstrauss and Tsukamoto
\cite[Lemma~4.2]{GutmanLindenstraussTsukamoto2016}.
For a tower of height $L$ with negative iterates and base $\Omega$, take
$V=h^{-(L-1)}\Omega$: its positive levels give the same tower in
reverse order.

\begin{theorem}[Niu's open tower theorem]
\label{thm:niu-open-tower}
For every integer $L\geq1$ and every $\varepsilon>0$, there is an open
set $V\subset X$ such that $V,hV,\ldots,h^{L-1}V$ are pairwise disjoint
and, writing $\mathcal T=\bigcup_{j=0}^{L-1}h^jV$,
\[
 \lim_{m\to\infty}\frac1m\sup_{x\in X}
   \sum_{s=0}^{m-1}\mathbf 1_{X\setminus\mathcal T}(h^sx)
 <\varepsilon.
\]
In particular, $\mu(\mathcal T)>1-\varepsilon$ for every
$\mu\in\mathcal M_h(X)$.
\end{theorem}

Niu's theorem assumes the marker property.  That hypothesis holds here:
there are no periodic points, since a finite orbit would be a nonempty
closed invariant subset of $X$.  Given a finite collection of nonzero
integers, freeness provides a nonempty open set disjoint from its
translates by those integers, and minimality makes all its translates
cover $X$.  This is the marker property.  Thus
\Cref{thm:niu-open-tower} requires no restriction on the covering
dimension of $X$.  The assertion about invariant measures follows by
integrating each finite orbit average before taking the limit.

For positive integers $k,H_0$ with $H_0\ge2$, put $e=kH_0$.
Whatever the offset $0\le i<k$, the $H_0-1$ length-$k$ intervals
\[
 [i+Dk,i+(D+1)k),\qquad 0\le D\le H_0-2,
\]
lie in $[0,e)$. The next lemma supplies a common open set $B$ for
their starting states and disjoint open regions $O_i$ on which the
offset $i$ can be used. A continuous weight $w$ vanishes outside
these regions and equals one at all but a uniformly small fraction
of the $e$-spaced starts along every orbit.

\begin{lemma}[Uniform block allocation]
\label{lem:uniform-block-allocation}
Let $k,H_0$ be positive integers with $H_0\geq2$, put $e=kH_0$, and let
$b_0>0$.  Suppose that $\gamma>0$ and $r\geq0$ satisfy
$\gamma/k>r$.  There exist an open set $B\subset X$, pairwise disjoint
open sets $O_0,\ldots,O_{k-1}\subset X$, a continuous function
$w:X\to[0,1]$, a compact set $K\subset B$, and an integer $J_0\geq1$
with the following properties:
\begin{enumerate}
\item $B,hB,\ldots,h^{k-1}B$ are pairwise disjoint, and
      $\gamma\mu(B)>r$ for every $\mu\in\mathcal M_h(X)$.
\item For $0\leq i<k$ and $0\leq D\leq H_0-2$,
      \begin{equation}
      \label{eq:allocated-block-starts}
        h^{i+Dk}O_i\subset B.
      \end{equation}
\item The support of $w$ is contained in $O=\bigcup_{i=0}^{k-1}O_i$,
      and
      \begin{equation}
      \label{eq:compact-selected-starts}
       h^{i+Dk}(\operatorname{supp}(w)\cap O_i)\subset K
       \quad(0\leq i<k,\ 0\leq D\leq H_0-2).
      \end{equation}
\item For every $x\in X$, every $v\in\mathbb Z$, and every $J\geq J_0$,
      \begin{equation}
      \label{eq:uniform-bad-block-count}
       \#\{0\leq q<J:w(h^{v+qe}x)<1\}\leq b_0J.
      \end{equation}
\end{enumerate}
In particular, the alignment label $i$ is locally constant on a
neighborhood of $\operatorname{supp}(w)$.
\end{lemma}

\begin{proof}
Choose $\eta>0$ so small that
\begin{equation}
\label{eq:allocation-error-budget}
 \frac\gamma k(1-\eta)>r,
 \qquad 3\eta<\frac{b_0}{4e},
 \qquad \eta<\frac18.
\end{equation}
Choose a multiple $L$ of $k$ with $L>e$ and $e/L<\eta$.
Apply \Cref{thm:niu-open-tower} with height $L$ and error $\eta$, and
denote its base by $V$.  Set
\[
 B=\bigcup_{a=0}^{L/k-1}h^{ak}V.
\]
Its first $k$ translates partition the tower.  Hence they are pairwise
disjoint, and
\[
 k\mu(B)=\mu\left(\bigcup_{s=0}^{L-1}h^sV\right)>1-\eta.
\]
The first requirement follows from
\Cref{eq:allocation-error-budget}.

For each integer $s$ with $0\leq s<L-e$, let $i(s)$ be the unique
integer in $\{0,\ldots,k-1\}$ satisfying $s+i(s)\equiv0\pmod{k}$,
and put
\[
 O_i=\bigcup_{\substack{0\leq s<L-e\\ i(s)=i}}h^sV.
\]
These are pairwise disjoint open sets.  For the indicated values of
$D$, the integer $s+i(s)+Dk$ is a multiple of $k$ between $0$ and
$L-1$, which proves \Cref{eq:allocated-block-starts}.  Also
\[
 \mu(O)\geq\mu\left(\bigcup_{s=0}^{L-1}h^sV\right)-e\mu(V)
 >1-\eta-\frac eL>1-2\eta.
\]

We next shrink $O$ uniformly over invariant measures.  Take continuous
functions $\varphi_j:X\to[0,1]$ increasing pointwise to
$\mathbf 1_O$, each with compact support in $O$.  Such functions are
obtained from increasing cutoffs of the distance to $X\setminus O$;
if $O=X$, use the constant function one.  For each invariant measure,
monotone convergence gives
$\int\varphi_j\,d\mu\to\mu(O)>1-2\eta$.
The sets
\[
 \left\{\mu\in\mathcal M_h(X):
                  \int\varphi_j\,d\mu>1-3\eta\right\}
\]
are increasing open sets covering the compact space
$\mathcal M_h(X)$.  One of them therefore equals that whole space.
Fix the corresponding function $\varphi$ and set $U=\{\varphi>0\}$.
The continuous function $\mu\mapsto\int\varphi\,d\mu$ attains its
minimum, which is still strictly greater than $1-3\eta$.
Since $\mu(U)\geq\int\varphi\,d\mu$, we obtain
$\overline U\subset O$ and
\begin{equation}
\label{eq:closed-exception-measure}
 C=X\setminus U\text{ is closed},\qquad
 \sup_{\mu\in\mathcal M_h(X)}\mu(C)<3\eta<\frac{b_0}{4e}.
\end{equation}
Choose $w:X\to[0,1]$ equal to one on $\overline U$ and with
$\operatorname{supp}(w)\subset O$.

For completeness, the uniform orbit estimate for the closed set $C$ is
\begin{equation}
\label{eq:closed-set-uniform-orbit-estimate}
 \limsup_{m\to\infty}\sup_{x\in X}\frac1m
       \sum_{s=0}^{m-1}\mathbf 1_C(h^sx)
 \leq\sup_{\mu\in\mathcal M_h(X)}\mu(C).
\end{equation}
Indeed, a sequence violating this estimate has a subsequence of its
orbit empirical measures converging weakly to an invariant probability
measure.  The closed-set inequality for weak convergence bounds the
limsup of their masses on $C$ by the mass of that limit measure, a
contradiction.  Thus, by
\Cref{eq:closed-exception-measure,eq:closed-set-uniform-orbit-estimate},
there is $m_0$ such that every orbit interval of length $m\geq m_0$
spends less than $b_0/(2e)$ of its times in $C$.

The set $\{w<1\}$ is contained in $C$.  For $eJ\geq m_0$ and any
$y\in X$,
\[
 \begin{split}
 \#\{0\leq q<J:w(h^{qe}y)<1\}
 &\leq\sum_{s=0}^{eJ-1}\mathbf 1_C(h^sy)\\
 &<\frac{b_0}{2e}\,eJ\leq b_0J.
 \end{split}
\]
Taking $y=h^vx$ proves \Cref{eq:uniform-bad-block-count} for a suitable
$J_0$.  This argument counts all integer times before passing to
$e$-spaced starts; it makes no assumption about invariant measures for
$h^e$.

Finally, each $\operatorname{supp}(w)\cap O_i$ is compact: the $O_i$
form a finite open partition of a neighborhood of
$\operatorname{supp}(w)$.  Let $K$ be the union of their images under
the finitely many maps in \Cref{eq:compact-selected-starts}.
It is a compact subset of $B$ by
\Cref{eq:allocated-block-starts}, as required.
\end{proof}

\subsection{A fixed linear map for changing coordinate supports}

The following elementary observation will be used twice.  It avoids
making a discontinuous choice of a projection for each support pattern.

\Needspace{6\baselineskip}
\begin{lemma}
\label{lem:finite-coordinate-compression}
Let $E$ be a finite index set and let $b,q$ be nonnegative integers.
There is a complex linear map
\[
 L:\bigoplus_{s\in E}\mathbb C^b\longrightarrow\mathbb C^{bq}
\]
whose restriction to every coordinate subspace supported on at most
$q$ indices is injective.
\end{lemma}

\begin{proof}
For each $\Lambda\subset E$ with $|\Lambda|\leq q$, failure of
injectivity on $\bigoplus_{s\in\Lambda}\mathbb C^b$ is a proper
determinantal algebraic condition on $L$: it is proper because a map
injective on that particular subspace plainly exists.  There are only
finitely many such conditions, so their union does not fill the vector
space of linear maps.  Any map outside their union works.  The cases
$b=0$ or $q=0$ are immediate.
\end{proof}

Only injectivity on each individual coordinate subspace is asserted.
Injectivity on the sum of all possible support subspaces would require
a much larger target and is not needed.

\subsection{The parametrized bundle embedding}

Fix $p\geq1$ and put $t=2p+1$.  Use the framed torus
$L\subset(-1/32,1/32)^t$ and the rank-$p$ bundle
$H\subset L\times\mathbb C^{2p}$ from
\Cref{thm:torus-compression}. Recall that
$Q:L\to M_{2p}(\mathbb C)$ is the orthogonal projection onto $H$.
In the notation of
\Cref{eq:sparse-slots}, $Z_{an}(t,n)$ consists of the parameters with at
most $an$ nonzero whole slots.

\begin{proposition}[Comparison gives a parametrized bundle embedding]
\label{prop:comparison-bundle}
Let $r$ be a comparison bound for $A$.  Fix positive integers
$p,g,k,H_0$ with $H_0\geq2$, numbers $0<a<1/16$ and $b_0>0$, and an
integer $R$ such that
\begin{equation}
\label{eq:comparison-bundle-budgets}
 \frac{RH_0}{pg(H_0-1)}+2(a+b_0)<2,
 \qquad \frac{R-pg}{k}>r.
\end{equation}
Let $F_1,\ldots,F_g:X\to[-1,1]^t$ be arbitrary continuous maps.
Choose $B,(O_i)_{i=0}^{k-1},w,K$ as in
\Cref{lem:uniform-block-allocation}, with $\gamma=R-pg$ and
$e=kH_0$.

For all sufficiently large integers $J$, put
\[
 N=eJ,\qquad n=J(H_0-1)g.
\]
There is a continuous map
$f_J:X\to[-1,1]^{tn}$, with slots indexed by
\[
 (q,D,j),\qquad
 0\leq q<J,\quad 0\leq D\leq H_0-2,\quad 1\leq j\leq g,
\]
such that
\begin{equation}
\label{eq:aligned-slot-agreement}
 (f_J)_{q,D,j}(x)=F_j(h^{qe+i+Dk}x)
 \quad\text{when }h^{qe}x\in O_i
       \text{ and }w(h^{qe}x)>0.
\end{equation}
The agreement holds, more strongly, on a neighborhood of each compact
set $h^{-qe}(\operatorname{supp}(w)\cap O_i)$.

Over the compact space
\begin{equation}
\label{eq:dynamical-torus-locus}
 \mathcal L_J=
 \{(x,z)\in X\times Z_{an}(t,n): f_J(x)-z\in L^n\},
\end{equation}
the direct sum of the $n$ pullbacks of $H$, by the slot maps
$(x,z)\mapsto(f_J)_\alpha(x)-z_\alpha$, embeds continuously in a
trivial complex bundle of rank strictly less than $2pn$.
The same allocation $B,(O_i),w,K$ works for every such $J$.
\end{proposition}

The two inequalities in \Cref{eq:comparison-bundle-budgets} control
different parts of the construction. The tracial comparison will use
source rank at most $pg$ and target rank $R$ on $B$. Since the tower
can make $\mu(B)$ uniformly as close to $1/k$ as needed, the excess
$(R-pg)/k>r$ supplies the required strict trace gap.
For the bundle embedding, each of the $J$ intervals contributes
$(H_0-1)g$ source slots, giving rank $pn=pg(H_0-1)J$.
The comparison witness will have a fixed propagation bound $M$;
the $k$-separation of its active output times will give a target of
rank $RH_0J+O(1)$. A second map for the sparse slots and the starts
where $w<1$ will use at most $2p(a+b_0)n+O(1)$ further dimensions.
The bounded terms are independent of $J$. Dividing these target
costs by $pn$ explains the first inequality: it leaves a strict margin
below the factor two required by \Cref{thm:torus-compression}.

\begin{proof}
We retain the allocation data throughout the proof.  In particular, all
cutoffs and the Cuntz witness below are fixed before $J$ is chosen.

\paragraph{Extension of the aligned functions.}
Put $W_i=\operatorname{supp}(w)\cap O_i$ and choose open sets $V_i$
with
\[
 W_i\subset V_i\subset\overline V_i\subset O_i;
\]
an empty $W_i$ may be assigned $V_i=\varnothing$.  For fixed $D,j$,
the maps $x\mapsto F_j(h^{i+Dk}x)$ on the disjoint closed sets
$\overline V_i$ combine to a continuous cube-valued map.  Coordinatewise
extension gives a continuous map
$\widehat F_{D,j}:X\to[-1,1]^t$ agreeing with it on their union.
For every $J$, set
\begin{equation}
\label{eq:extended-aligned-functions}
 (f_J)_{q,D,j}(x)=\widehat F_{D,j}(h^{qe}x).
\end{equation}
This proves the asserted agreement, including the neighborhood
assertion.  The extension maps $\widehat F_{D,j}$ are independent of
$J$.

\paragraph{Positive matrix functions.}
We now encode the $g$ torus bundles by positive elements of the crossed
product. Their tracial rank gap will supply a single witness for all the
aligned orbit blocks.
Take a tubular neighborhood $\mathcal V$ of $L$ inside $(-1,1)^t$,
with retraction $\pi_L:\mathcal V\to L$, and write
$\widetilde Q=Q\circ\pi_L$.  Choose a continuous function
$\psi:[-1,1]^t\to[0,1]$ with compact support in $\mathcal V$ and
equal to one on a neighborhood of $L$.  The matrix function
$\psi\widetilde Q$ extends by zero outside $\mathcal V$.
By compactness there is $\delta>0$ such that, for $y\in L$ and
$v\in[-1,1]^t$,
\begin{equation}
\label{eq:nearby-bundle-projections}
 \|v-y\|_\infty<\delta
 \quad\Longrightarrow\quad
 \psi(v)=1,\quad v\in\mathcal V,\quad
 \|\widetilde Q(v)-Q(y)\|<\frac12.
\end{equation}
Consequently $\widetilde Q(v)$ maps $H_y$ injectively onto
$\operatorname{ran}\widetilde Q(v)$ in this situation.

Choose continuous scalar functions $\rho_0,\rho_1:X\to[0,1]$ such that
\[
 \rho_0|_K=1,\qquad \operatorname{supp}(\rho_0)\subset B,
 \qquad \{\rho_1>0\}=B.
\]
The compact containment $K\subset B$ gives $\rho_0$, and a cutoff of
the distance to $X\setminus B$ gives $\rho_1$; if $B=X$, take
$\rho_1=1$.  Define positive matrix functions
$P_0\in M_{2pg}(C(X))_+$ and $P_1\in M_R(C(X))_+$ by
\begin{equation}
\label{eq:comparison-positive-functions}
 \begin{split}
 P_0(x)&=\rho_0(x)\bigoplus_{j=1}^g
              \psi(F_j(x))\widetilde Q(F_j(x)),\\
 P_1(x)&=\rho_1(x)1_R.
 \end{split}
\end{equation}
The summands of $P_0$ are understood to be zero where $F_j(x)$ lies
outside $\mathcal V$.

If $\tau$ is a normalized trace on $A$, its restriction to $C(X)$ is
integration against some $\mu\in\mathcal M_h(X)$.  Functional calculus
and bounded convergence give, for every positive matrix function $P$,
\[
 d_\tau(P)=\int_X\operatorname{rank}(P(x))\,d\mu(x).
\]
The unnormalized matrix trace is essential in this formula.  The rank
of $P_0(x)$ is at most $pg$ on $B$ and is zero off $B$, whereas $P_1$
has rank $R$ precisely on $B$.  Therefore
\[
 d_\tau(P_0)+r\leq pg\mu(B)+r<R\mu(B)=d_\tau(P_1).
\]
Here the strict inequality is supplied by
\Cref{lem:uniform-block-allocation}.  Since $r$ is a comparison bound, it follows
that
\begin{equation}
\label{eq:allocated-cuntz-comparison}
 P_0\precsim_A P_1.
\end{equation}

\paragraph{A finite-propagation witness and simultaneous injection.}
Finite-support compression witnesses and bundle obstructions also enter
the lower-bound arguments of Hirshberg--Phillips
\cite[proofs of Theorems~3.3 and~4.5 in the author preprint]{HirshbergPhillips2022}.
Here we retain the sparse parameter throughout and require simultaneous
injectivity on the entire sum of the selected fibers.
Taking the appropriate rectangular corner of a Cuntz witness for
\Cref{eq:allocated-cuntz-comparison}, and then approximating its entries
by finite Fourier sums, gives
$T_0\in M_{R,2pg}(A)$ such that
\begin{equation}
\label{eq:finite-propagation-witness}
 \|T_0^*P_1T_0-P_0\|<\frac12,
 \qquad
 T_0=\sum_{\ell=-M}^{M} A_\ell u^\ell
\end{equation}
for some fixed $M\geq0$ and continuous matrix functions $A_\ell$.
For example, first choose a witness with error less than $1/4$; norm
continuity of $T\mapsto T^*P_1T$ permits the Fourier approximation while
keeping the error below $1/2$.

In the orbit representation at $x$, a function $a\in C(X)$ acts by
\[
 (\pi_x(a)\xi)(s)=a(h^sx)\xi(s),\qquad s\in\mathbb Z,
\]
and the implementing unitary acts by a shift, with direction determined
by the crossed-product convention.  Thus $T_0$ has propagation at most
$M$.  Let
\[
 I_J=\{-M,-M+1,\ldots,N-1+M\}.
\]
Restricting the input to times $0,\ldots,N-1$ and the output to $I_J$
gives a continuous finite matrix
\[
 W_J(x):\bigoplus_{s=0}^{N-1}\mathbb C^{2pg}
          \longrightarrow\bigoplus_{s\in I_J}\mathbb C^R
\]
representing $\pi_x(P_1)^{1/2}\pi_x(T_0)$ on that input window.
There is no output outside $I_J$.  If an input vector $\xi$ in the
window satisfies $\pi_x(P_0)\xi=\xi$, then
\begin{equation}
\label{eq:whole-projection-sum-injection}
 \|W_J(x)\xi\|^2
 =\langle\pi_x(T_0^*P_1T_0)\xi,\xi\rangle
 \geq\frac12\|\xi\|^2.
\end{equation}
This holds on the whole sum of any selected genuine projection ranges,
including sums involving different times and different bundle blocks.

\paragraph{The rank cost of separated output rows.}
The map $W_J(x)$ initially has $(N+2M)R$ output coordinates.
We now bound the target rank needed to retain the injectivity in
\Cref{eq:whole-projection-sum-injection}, using the spacing of the
rows that can be nonzero.
The output row at time $s$ vanishes unless $h^sx\in B$.  These active
times are separated by at least $k$, because
$B,hB,\ldots,h^{k-1}B$ are disjoint.  There are at most
\[
 q_J=\left\lceil\frac{N+2M}{k}\right\rceil
       =H_0J+c_M,\qquad c_M=\left\lceil\frac{2M}{k}\right\rceil,
\]
of them in $I_J$.  By \Cref{lem:finite-coordinate-compression}, choose
one fixed linear map
\[
 C_J:\bigoplus_{s\in I_J}\mathbb C^R\longrightarrow\mathbb C^{Rq_J}
\]
injective on every coordinate subspace supported at at most $q_J$
times.  Hence $C_J$ is injective on the range of $W_J(x)$ for every
$x$.  The matrix $C_J$ is chosen once for this finite window and does
not depend on $x$ or on the sparse parameter.

\paragraph{The first bundle map.}
Let $\alpha=(q,D,j)$ index a slot.  On $\mathcal L_J$, write
\[
 y_\alpha=(f_J)_\alpha(x)-z_\alpha\in L,
 \qquad
 E_{(x,z)}=\bigoplus_\alpha H_{y_\alpha}
                  \subset\bigoplus_\alpha\mathbb C^{2p}.
\]
This describes the rank-$pn$ bundle $E$ in the proposition.
Choose a continuous $\chi:[-2,2]^t\to[0,1]$ with $\chi(0)=1$ and
support contained in $\{z:\|z\|_\infty<\delta\}$.  Set
\begin{equation}
\label{eq:slot-weights}
 \theta_\alpha(x,z)=w(h^{qe}x)\chi(z_\alpha).
\end{equation}

Where $\theta_\alpha>0$, there is a unique $i$ with $h^{qe}x\in O_i$.
Put $s_\alpha=qe+i+Dk$.  By
\Cref{eq:aligned-slot-agreement,eq:nearby-bundle-projections},
\[
 \widetilde Q(F_j(h^{s_\alpha}x)):
 H_{y_\alpha}\longrightarrow
        \operatorname{ran}\widetilde Q(F_j(h^{s_\alpha}x))
\]
is injective.  Furthermore, $h^{s_\alpha}x\in K$, so the indicated
range is a genuine projection range of the $j$th block of
$P_0(h^{s_\alpha}x)$.

Define a bundle map
\[
 \Xi:E\longrightarrow
      \mathcal L_J\times\bigoplus_{s=0}^{N-1}\mathbb C^{2pg}
\]
by sending $v_\alpha\in H_{y_\alpha}$, where
$\theta_\alpha>0$, to
\[
 \theta_\alpha\widetilde Q(F_j(h^{s_\alpha}x))v_\alpha
\]
in the $j$th block of the row at time $s_\alpha$, and by setting that
contribution to zero where $\theta_\alpha=0$.
All such times lie in the $q$th interval $[qe,(q+1)e)$.
Within an interval the pairs $(i+Dk,j)$ are distinct, and intervals
are disjoint.  Thus different active slots occupy different
time/block pairs.

This map is continuous on the entire space $\mathcal L_J$.
Indeed, on a region with positive weight the label is locally constant
and the formula is continuous.  The projection factor has norm at most
one, so the norm of its weighted contribution is at most
$\theta_\alpha\|v_\alpha\|$ and tends to zero wherever the weight
vanishes.  This also handles a change of label, which can occur only
outside the positive-weight region.  More formally, one can use the
matrix $\widetilde Q(F_j(h^{s_\alpha}x))Q(y_\alpha)$ on the ambient
$\mathbb C^{2p}$ and extend its weighted entries by zero; the same
bound proves their continuity.  No nearby-projection map is needed at
zero weight.

The image of $\Xi_{(x,z)}$ lies in the subspace on which
$\pi_x(P_0)$ is the identity.  The first bundle map is therefore
\[
 \Phi_1(x,z)=C_JW_J(x)\Xi_{(x,z)}:
 E_{(x,z)}\longrightarrow\mathbb C^{Rq_J}.
\]
By \Cref{eq:whole-projection-sum-injection}, the defining property of
$C_J$, and the distinct time/block placements,
\begin{equation}
\label{eq:first-map-kernel}
 \ker\Phi_1(x,z)=
       \bigoplus_{\alpha:\theta_\alpha(x,z)=0} H_{y_\alpha}.
\end{equation}

\paragraph{The exceptional summands.}
The first map now controls every positive-weight summand. To complete
one embedding over the full parameter locus, we give its kernel a second
map whose target rank is charged only to the sparse slots and the
infrequent starts at which $w<1$.
Using the given embeddings $H_y\subset\mathbb C^{2p}$, define
\[
 \Psi_{(x,z)}:E_{(x,z)}\longrightarrow
                     \bigoplus_\alpha\mathbb C^{2p},
 \qquad
 (v_\alpha)_\alpha\longmapsto
                 ((1-\theta_\alpha)v_\alpha)_\alpha.
\]
If $1-\theta_\alpha\ne0$, then either $z_\alpha\ne0$ or
$w(h^{qe}x)<1$: when both $z_\alpha=0$ and $w(h^{qe}x)=1$,
\Cref{eq:slot-weights} gives $\theta_\alpha=1$.
The sparse condition bounds the first kind by $an$ slots.  For
$J\geq J_0$, \Cref{eq:uniform-bad-block-count} bounds the second kind
by
\[
 b_0J(H_0-1)g=b_0n
\]
slots.  The range of $\Psi_{(x,z)}$ is therefore supported on at most
$(a+b_0)n$ coordinate groups, each of dimension $2p$.

Put $s_J=\lceil(a+b_0)n\rceil$.  Another application of
\Cref{lem:finite-coordinate-compression} gives a fixed linear map
\[
 D_J:\bigoplus_\alpha\mathbb C^{2p}\longrightarrow
                         \mathbb C^{2ps_J}
\]
injective on all the relevant coordinate subspaces.  Set
$\Phi_2(x,z)=D_J\Psi_{(x,z)}$.
The direct sum $\Phi_1\oplus\Phi_2$ is fiberwise injective.  In fact,
\Cref{eq:first-map-kernel} says that a vector killed by $\Phi_1$ has
components only where $\theta_\alpha=0$.  On those components
$1-\theta_\alpha=1$.  Injectivity of $D_J$ on the range support of
$\Psi$ shows that a vector also killed by $\Phi_2$ must be zero.

Both maps are continuous in $(x,z)$, including when the support of $z$
changes.  They thus give a continuous bundle embedding into the
trivial bundle of rank
\begin{equation}
\label{eq:comparison-bundle-target-rank}
 K_J=Rq_J+2ps_J.
\end{equation}
One may also verify the bundle assertion locally by representing the
map in a trivialization of $E$: its Gram matrix is positive definite,
so its image is a continuous rank-$pn$ subbundle.  The base
$\mathcal L_J$ is compact.

\paragraph{The strict rank estimate.}
Let
\[
 \zeta=2-\frac{RH_0}{pg(H_0-1)}-2(a+b_0)>0.
\]
Since $n=J(H_0-1)g$, \Cref{eq:comparison-bundle-target-rank} gives
\[
 \frac{K_J}{pn}
 \leq \frac{RH_0}{pg(H_0-1)}+2(a+b_0)
            +\frac{Rc_M+2p}{pJ(H_0-1)g}
 =2-\zeta+\frac{Rc_M+2p}{pJ(H_0-1)g}.
\]
For all sufficiently large $J\geq J_0$ this is strictly less than
two.  Hence $K_J<2pn$, which proves the proposition.
\end{proof}
